%\RequirePackage[l2tabu, orthodox]{nag}%nag complains because beamer titlepage uses a centering environment
%\WarningFilter{nag}{There is no environment ``centering'' }%this leaves nag complaining about “1 complaints in total”.
\RequirePackage{silence}
\WarningFilter{ifpdf}{Someone has redefined \pdfoutput}
\WarningFilter{fmtcount}{\ordinal already defined use \FCordinal instead}
\documentclass[english]{beamer}
%INSTALL

%avoids a warning
\usepackage[log-declarations=false]{xparse}
\usepackage{fontspec} %font selecting commands
\usepackage{xunicode}
\usepackage{ucharclasses}
%warn about missing characters
\tracinglostchars=2

%REDAC
\usepackage{booktabs}
\usepackage{calc}
\usepackage{tabularx}

\usepackage{etoolbox} %for addtocmd, newtoggle commands
\newtoggle{LCpres}
\toggletrue{LCpres}

\usepackage{mathtools} %load this before babel!
	\mathtoolsset{showonlyrefs,showmanualtags}

%\usepackage[french, english]{babel}%options should probably be on the document level
\usepackage{babel}
%suppresses the warning about frenchb not modifying the captions (“—” to “:” in “Figure 1 – Legend”).

%\usepackage[super]{nth}%better use fmtcount! (loaded by datetime anyway; see below about pbl with warnings and package silence)
\usepackage{listings} %typeset source code listings
	\lstset{language=XML,tabsize=2,captionpos=b,basicstyle=\NoAutoSpacing}%NoAutoSpacing avoids space before colon or ?}%,literate={"}{{\tt"}}1, keywordstyle=\fontspec{Latin Modern Mono Light}\textbf, emph={String, PreparedStatement}, emphstyle=\fontspec{Latin Modern Mono Light}\textbf, language=Java, basicstyle=\small\NoAutoSpacing\ttfamily, aboveskip=0pt, belowskip=0pt, showstringspaces=false
\usepackage[nolist,smaller,printonlyused]{acronym}%,smaller option produces warnings from relsize in some cases, it seems.% Note silence and acronym and hyperref make (xe)latex crash when ac used in section (http://tex.stackexchange.com/questions/103483/strange-packages-interaction-acronyms-silence-hyperref), rather use \section{\texorpdfstring{\acs{UE}}{UE}}.
\usepackage{fmtcount}
\usepackage[nodayofweek]{datetime}%must be loaded after the babel package. However, loading it after {nth} generates a warning from fmtcount about ordinal being already defined. Better load it before nth? (then we can remove the silence package which creates possible crashes, see above.) Or remove nth?
%\usepackage{xspace}%do we need this?
\nottoggle{LCpres}{
	\usepackage[textsize=small]{todonotes}
}
\iftoggle{LCpres}{
	%remove pdfusetitle (implied by beamer)
	\usepackage{hyperref}
}{
% option pdfusetitle must be introduced here, not in hypersetup.
	\usepackage[pdfusetitle]{hyperref}
}
%breaklinks makes links on multiple lines into different PDF links to the same target.
%colorlinks (false): Colors the text of links and anchors. The colors chosen depend on the the type of link. In spite of colored boxes, the colored text remains when printing.
%linkcolor=black: this leaves other links in colors, e.g. refs in green, don't print well.
%pdfborder (0 0 1, set to 0 0 0 if colorlinks): width of PDF link border
%hidelinks
\hypersetup{breaklinks, bookmarksopen}
%add hyperfigures=true in hypersetup (already defined in article mode)
\iftoggle{LCpres}{
	\hypersetup{hyperfigures}
}

%in Beamer, sets url colored links but does not change the rest of the colors (http://tex.stackexchange.com/questions/13423/how-to-change-the-color-of-href-links-for-real)
%\hypersetup{breaklinks,bookmarksopen,colorlinks=true,urlcolor=blue,linkcolor=,hyperfigures=true}
% hyperref doc says: Package bookmark replaces hyperref’s bookmark organization by a new algorithm (...) Therefore I recommend using this package.
\usepackage{bookmark}

% center floats by default, but do not use with float
% \usepackage{floatrow}
% \makeatletter
% \g@addto@macro\@floatboxreset\centering
% \makeatother
\nottoggle{LCpres}{
	\usepackage{enumitem} %follow enumerate by a string saying how to display enumeration
}
\usepackage{ragged2e} %new com­mands \Cen­ter­ing, \RaggedLeft, and \RaggedRight and new en­vi­ron­ments Cen­ter, FlushLeft, and FlushRight, which set ragged text and are eas­ily con­fig­urable to al­low hy­phen­ation (the cor­re­spond­ing com­mands in LaTeX, all of whose names are lower-case, pre­vent hy­phen­ation al­to­gether). 
\usepackage{siunitx} %[expproduct=tighttimes, decimalsymbol=comma] ou (plus récent ?) [round-mode=figures, round-precision=2, scientific-notation = engineering]
\sisetup{detect-all, locale = FR, strict}% to detect e.g. when in math mode (use a math font) - check whether this makes sense with strict
\usepackage{braket} %for \Set
\usepackage{natbib}
\usepackage{doi}

\usepackage{amsmath,amsthm}
% \usepackage{amsfonts} %not required?
% \usepackage{dsfont} %for what?
%unicode-math overwrites the following commands from the mathtools package: \dblcolon, \coloneqq, \Coloneqq, \eqqcolon. Using the other colon-like commands from mathtools will lead to inconsistencies. Plus, Using \overbracket and \underbracke from mathtools package. Use \Uoverbracket and \Uunderbracke for original unicode-math definition.
%use exclusively \mathbf and choose math bold style below.
\usepackage[warnings-off={mathtools-colon, mathtools-overbracket}, bold-style=ISO]{unicode-math}

\defaultfontfeatures{
	Fractions=On,
	Mapping=tex-text% to turn "--" into dashes, useful for bibtex%%
}
%\defaultfontfeatures[\rmfamily, \sffamily]{
%\defaultfontfeatures[\rmfamily]{
%	Fractions=On,
%	Mapping=% to leave " alone (disable the default mapping tex-text; requires loading the font afterwards?)
%}
\newfontfamily\xitsfamily{XITS}
\newfontfamily\texgyretermesfamily{TeX Gyre Termes}
\newfontfamily\texgyreherosfamily{TeX Gyre Heros}
\newfontfamily\lmfamily{Latin Modern Roman}
\iftoggle{LCpres}{
	\setmainfont{Latin Modern Roman}
	\setsansfont{Latin Modern Sans}
	\setmonofont{Latin Modern Mono}
}{
	\setmainfont{TeX Gyre Termes}
	\setsansfont{TeX Gyre Heros}
	\setmonofont{TeX Gyre Cursor}
}
\defaultfontfeatures{}%disable default font features to avoid warnings with math fonts.
\setmathfont{XITS Math}
\setmathfont[range={\mathcal,\mathbfcal},StylisticSet=1]{XITS Math}
%silently delays restoring normal font when using e.g. “≠)” or “≠ other word”.
%\setTransitionTo{Arrows}{\xitsfamily}
%\setTransitionTo{MathematicalOperators}{\xitsfamily}
%this solution explicitly fails on “≠)” and works for “≠ other word.”
\setTransitionsFor{MathematicalOperators}{\begingroup\xitsfamily}{\endgroup}
\setTransitionsFor{Arrows}{\begingroup\xitsfamily}{\endgroup}
%texgyretermes does not have ✓, XITS does.
\setTransitionsFor{Dingbats}{\begingroup\xitsfamily}{\endgroup}
%this also works, but it is more complex
%\def\ResetTransitionTo#1{%
%  \XeTeXinterchartoks 255 \csname#1Class\endcsname{\relax}}
%\setTransitionsFor{MathematicalOperators}
%  {\begingroup\ResetTransitionTo{MathematicalOperators}\xitsfamily}
%  {\endgroup}

\usepackage{cleveref}% cleveref should go "laster" than hyperref
%GRAPHICS
\usepackage{pgf}
\usepackage{pgfplots}
	\usetikzlibrary{babel, matrix, fit, plotmarks, calc, trees, shapes.geometric, positioning, plothandlers, arrows, shapes.multipart}
\pgfplotsset{compat=1.11}
\usepackage{graphicx}

\graphicspath{{graphics/},{graphics-dm/}}
\DeclareGraphicsExtensions{.pdf}
\LetLtxMacro\SavedIncludeGraphics\includegraphics
\AtBeginDocument{
	\def\includegraphics#1#{% #1 catches optional stuff (star/opt. arg.)
		\IncludeGraphicsAux{#1}%
	}%
}
\newcommand*{\IncludeGraphicsAux}[2]{%
	\XeTeXLinkBox{%
		\SavedIncludeGraphics#1{#2}%
	}%
}%

%HACKING
\usepackage{printlen}
\uselengthunit{mm}
% 	\newlength{\templ}% or LenTemp?
% 	\setlength{\templ}{6 pt}
% 	\printlength{\templ}
\usepackage{scrhack}% load at end. Corrects a bug in float package, which is outdated but might be used by other packages
\usepackage{xltxtra} %somebody said that this is loaded by fontspec, but does not seem correct: if not loaded explicitly, does not appear in the log and \showhyphens is not corrected.

%Beamer-specific
%do not remove babel, which beamer uses (beamer uses the \translate command for the appendix); but french can be removed.
\iftoggle{LCpres}{
	\usepackage{appendixnumberbeamer}
	\setbeamertemplate{navigation symbols}{} 
	\usepackage{preamble/beamerthemeParisFrance}
	\usefonttheme{professionalfonts}
	\setcounter{tocdepth}{10}
	%From: http://tex.stackexchange.com/questions/168057/beamer-with-xelatex-on-texlive2013-enumerate-numbers-in-black
%I don’t think it’s useful to submit this as a bug: nothing has been solved since March, 2015. See: https://bitbucket.org/rivanvx/beamer/issues?status=resolved.

\setbeamertemplate{enumerate item}
{
  \begin{pgfpicture}{-1ex}{-0.65ex}{1ex}{1ex}
    \usebeamercolor[fg]{item projected}
    {\pgftransformscale{1.75}\pgftext{\normalsize\pgfuseshading{bigsphere}}}
    {\pgftransformshift{\pgfpoint{0pt}{0.5pt}}
      \pgftext{\usebeamercolor[fg]{item projected}\usebeamerfont*{item projected}\insertenumlabel}}
  \end{pgfpicture}%
}

\setbeamertemplate{enumerate subitem}
{
  \begin{pgfpicture}{-1ex}{-0.55ex}{1ex}{1ex}
    \usebeamercolor[fg]{subitem projected}
    {\pgftransformscale{1.4}\pgftext{\normalsize\pgfuseshading{bigsphere}}}
    \pgftext{%
      \usebeamercolor[fg]{subitem projected}%
      \usebeamerfont*{subitem projected}%
      \insertsubenumlabel}
  \end{pgfpicture}%
}

\setbeamertemplate{enumerate subsubitem}
{
  \begin{pgfpicture}{-1ex}{-0.55ex}{1ex}{1ex}
    \usebeamercolor[fg]{subsubitem projected}
    {\pgftransformscale{1.4}\pgftext{\normalsize\pgfuseshading{bigsphere}}}
    \pgftext{%
      \usebeamercolor[fg]{subsubitem projected}%
      \usebeamerfont*{subitem projected}%
      \insertsubsubenumlabel}
  \end{pgfpicture}%
}


}
% \newcommand{\citep}{\cite}%Better: leave natbib.
% \setbeamersize{text margin left=0.1cm, text margin right=0.1cm} 
% \usetheme{BrusselsBelgium}%no, replace with paris
%\usetheme{ParisFrance}, no, usepackage better!
% Tex Gyre takes too much space, replace with Latin Modern Roman / Sans / Mono.
% Difference when loading explicitly Latin Modern Sans (compared to not using \setsansfont at all):
% the font LMSans17-Regular appears in the document;
% the title of the slides appears differently;
% it does not say (in the log file):
% > LaTeX Font Info:    Font shape `EU1/lmss/m/it' in size <10.95> not available
% > (Font)              Font shape `EU1/lmss/m/sl' tried instead on input line 85.
% > LaTeX Font Info:    Try loading font information for EU1+lmtt on input line 85.

%tikzposter-specific
%remove \usepackage{ragged2e}: causes 1=1 to be printed in the middle of the poster. (Anyway prints a warning about those characters being missing.)
%put [french, english] options next to \usepackage{babel} to avoid warning of 

\newcommand{\R}{ℝ}
\newcommand{\N}{ℕ}
\newcommand{\Z}{ℤ}
\newcommand{\card}[1]{\lvert{#1}\rvert}
\newcommand{\powerset}[1]{\mathscr{P}(#1)}%\mathscr rather than \mathcal: scr is rounder than cal (at least in XITS Math).
%powerset without zero
\newcommand{\powersetz}[1]{\mathscr{P}_\emptyset(#1)}
\newcommand{\suchthat}{\;\ifnum\currentgrouptype=16 \middle\fi|\;}
%\newcommand{\Rplus}{\reels^+\xspace}

\AtBeginDocument{%
	\renewcommand{\epsilon}{\varepsilon}
% we want straight form of \phi for mathematics, as recommended in UTR #25: Unicode support for mathematics.
%	\renewcommand{\phi}{\varphi}
}

% with amssymb, but I don’t want to use amssymb just for that.
% \newcommand{\restr}[2]{{#1}_{\restriction #2}}
%\newcommand{\restr}[2]{{#1\upharpoonright}_{#2}}
\newcommand{\restr}[2]{{#1|}_{#2}}%sometimes typed out incorrectly within \set.
%\newcommand{\restr}[2]{{#1}_{\vert #2}}%\vert errors when used within \Set and is typed out incorrectly within \set.
\DeclareMathOperator*{\argmax}{arg\,max}
\DeclareMathOperator*{\argmin}{arg\,min}


\newcommand{\allalts}{\mathcal{A}}
\newcommand{\alts}{A}
\newcommand{\allrules}{\mathcal{F}}

%the set of all profiles
\newcommand{\profiles}{\mathbfcal{R}}
%a rank-profile
\newcommand{\rprof}{\mathbf{x}}
% \newcommand{\rprof}{\text{\boldmath $x$}
\newcommand{\rankprofilesrvs}{\mathcal{X}}

% constraints over ranks
\newcommand{\prrefceq}{≥^C}
\newcommand{\prrefc}{>^C}
\newcommand{\prrefcsim}{\sim^C}

\newcommand{\wos}{\mathcal{W}}
\newcommand{\prs}{\mathcal{Z}}
%preorder induced by consistent C
\newcommand{\prCeq}{\ppglobeq^C}%\framebox{C}
\newcommand{\prC}{\ppglob^C}
\newcommand{\halfvoters}{\frac{n}{2}}
% dominates strictly
% \newcommand{\doms}{\RHD}
\newcommand{\doms}{\mathrel{\blacktriangleright}}
% \newcommand{\rvspace}{\mathbb{X}}
\newcommand{\rvspace}{\left[1, m\right]^N}
\newcommand{\cspace}{\mathcal{C}}
% pref over ranks (individually)
\newcommand{\prrefeq}{\succcurlyeq}
\newcommand{\prreff}{\succ}
% global weak-order over ranks
\newcommand{\pglobeq}{⪰}
%that’s not really what I want, but \not\succsim crashes xelatex
\def\not\pglobeq{\nsucceq}
%\newcommand{\pglob}{⪲}
\newcommand{\pglob}{⪲}
% global (partial) preorder over ranks
\newcommand{\ppglobeq}{\succsim}
\newcommand{\ppglobsim}{\sim}
\newcommand{\ppglob}{\succnsim}
%\newcommand{\ppglob}{\succ}
% \def\not\ppglobeq{\nsucceq}
%U+22E1
% \def\not\ppglobeq{⋡}
%U+2275
\def\not\ppglobeq{≵}

\newcommand{\Vpr}{V^{\,\prime}}

%Voting and MCDA
%\newcommand{\allalts}{\mathscr{A}}

%Voting
%rank-profile voting rule
\newcommand{\frp}{f_\text{r-p}}
\newcommand{\feasalts}{F}
\newcommand{\allvoters}{\mathscr{N}}
\newcommand{\voters}{N}
\newcommand{\allsystems}{\mathcal{G}}
\newcommand{\prof}{\mathbf{R}}
\newcommand{\allprofs}{\mathbfcal{R}}
\newcommand{\linors}{\mathcal{L}(\alts)}

\newcommand{\pbasic}[1]{\prof^{#1}_\epsilon}
\newcommand{\pelem}[1]{\prof^{#1}_e}
\newcommand{\pcycl}[1]{\prof^{#1}_c}
\newcommand{\pcycllong}[1]{\prof^{#1}_{cl}}
\newcommand{\pinv}[1]{\overline{\prof_{#1}}}
\newcommand{\dmap}{{\xitsfamily δ}}
%powerset without zero
\newcommand{\powersetz}[1]{\mathcal{P}_\emptyset(#1)}

%logic atom
%⟼ (long)
\DeclareDocumentCommand{\lato}{ O{\prof} O{\alts} }{[#1 \!⟼\! #2]}
%logic atom in
%↝, \stackrel{\in}{\mapsto}, ➲, ⥹
\newcommand{\tightoverset}[2]{%
  \mathop{#2}\limits^{\vbox to -.5ex{\kern-0.9ex\hbox{$#1$}\vss}}}
\DeclareDocumentCommand{\latoin}{ O{\prof} O{\alpha} }{[#1 \tightoverset{\in}{⟼} #2]}
\newcommand{\alllang}{\mathcal{L}}
\newcommand{\ltru}{\texttt{T}}
\newcommand{\lfal}{\texttt{F}}
\newcommand{\laxiom}[1]{{\texgyreherosfamily{\textsc{#1}}}}

%ARG TH
\newcommand{\AF}{\mathcal{AF}}
\newcommand{\labelling}{\mathcal{L}}
\newcommand{\labin}{\textbf{in}\xspace}
\newcommand{\labout}{\textbf{out}}
\newcommand{\labund}{\textbf{undec}\xspace}
\newcommand{\nonemptyor}[2]{\ifthenelse{\equal{#1}{}}{#2}{#1}}
\newcommand{\gextlab}[2][]{
	\labelling{\mathcal{GE}}_{(#2, \nonemptyor{#1}{\ibeatsr{#2}})}
}
\newcommand{\allargs}{A^*}
\newcommand{\args}{A}
\newcommand{\ar}{a}
\newcommand{\ext}{\mathcal{E}}

%MCDA+Arg
\newcommand{\dm}{d}
\newcommand{\ileadsto}{\rightcurvedarrow}
\newcommand{\mleadsto}[1][\eta]{\rightcurvedarrow_{#1}}
\newcommand{\ibeats}{\vartriangleright}
\newcommand{\mbeats}[1][\eta]{\vartriangleright_{#1}}

%MISC
\newcommand{\lequiv}{\Vvdash}
\newcommand{\weightst}{W^{\,t}}

%MCDA classical
\newcommand{\crits}{\mathcal{J}}

%Sorting
\newcommand{\cats}{\mathcal{C}}
\newcommand{\catssubsets}{2^\cats}
\newcommand{\catgg}{\vartriangleright}
\newcommand{\catll}{\vartriangleleft}
\newcommand{\catleq}{\trianglelefteq}
\newcommand{\catgeq}{\trianglerighteq}
\newcommand{\alttoc}[2][x]{(#1 \xrightarrow{} #2)}
\newcommand{\alttocat}[3]{(#2 \xrightarrow{#1} #3)}
\newcommand{\alttoI}{(x \xrightarrow{} \left[\underline{C_x}, \overline{C_x}\right])}
\newcommand{\alttocatdm}[3][t]{\left(#2 \thinspace \raisebox{-3pt}{$\xrightarrow{#1}$}\thinspace #3\right)}
\newcommand{\alttocatatleast}[2]{\left(#1 \thinspace \raisebox{-3pt}{$\xrightarrow[]{≥}$}\thinspace #2\right)}
\newcommand{\alttocatatmost}[2]{\left(#1 \thinspace \raisebox{-3pt}{$\xrightarrow[]{≤}$}\thinspace #2\right)}

\definecolor{darkgreen}{rgb}{0,0.6,0}
\newcommand{\commentOC}[1]{{\small\color{blue}{\selectlanguage{french}$\big[$OC: #1$\big]$}}}
%\newcommand{\commentOC}[1]{{\selectlanguage{french}{\todo{OC : #1}}}}
%Or: \todo[color=green!40]
\newcommand{\innote}[1]{{\scriptsize{#1}}}

%this probably requires outdated float package, see doc KomaScript for an alternative.
% \newfloat{program}{t}{lop}
% \floatname{program}{PM}

%definition, theorem, lemma, example environments, qed trickery are only needed in article mode (not Beamer)
\nottoggle{LCpres}{
%style is plain by default (italic text)
	\newtheorem{definition}{Definition}
	\newtheorem{theorem}{Theorem}
%no italic: expected.
%http://tex.stackexchange.com/questions/144653/italicizing-of-amsthm-package
	\newtheorem{lemma}{Lemma}
%\crefname{axiom}{axiom}{axioms}%might be needed for workaround bug in cref when defining new theorems?

%\ifdefined\theorem\else
%\newtheorem{theorem}{\iflanguage{english}{Theorem}{Théorème}}
%\fi

\theoremstyle{remark}
	\newtheorem{examplex}{Example}
	\newtheorem{remarkx}{Remark}

%trickery allowing use of \qedhere and the like.
\newenvironment{example}{
	\pushQED{\qed}\renewcommand{\qedsymbol}{$\triangle$}\examplex
}{
	\popQED\endexamplex
}
\newenvironment{remark}{
	\pushQED{\qed}\renewcommand{\qedsymbol}{$\triangle$}\remarkx
}{
	\popQED\endremarkx
}
}{
}
\crefname{examplex}{example}{examples}% I wonder why this is unnecessary in case of singular

%which line breaks are chosen: accept worse lines, therefore reducing risk of overfull lines. Default = 200
\tolerance=2000
%accept overfull hbox up to...
\hfuzz=2cm
%reduces verbosity about the bad line breaks
\hbadness 5000
%reduces verbosity about the underful vboxes
\vbadness=1300
%sloppy sets tolerance to 9999
\apptocmd{\sloppy}{\hbadness 10000\relax}{}{}

\bibliographystyle{abbrvnat}
%or \bibliographystyle{apalike} for presentations?

%doi package uses old-style dx.doi url, see 3.8 DOI system Proxy Server technical details, “Users may resolve DOI names that are structured to use the DOI system Proxy Server (http://doi.org (preferred) or http://dx.doi.org).”, https://www.doi.org/doi_handbook/3_Resolution.html
\makeatletter
\patchcmd{\@doi}{dx.doi.org}{doi.org}{}{}
\makeatother

% WRITING
%\newcommand{\ie}{i.e.\@\xspace}%to try
%\newcommand{\eg}{e.g.\@\xspace}
%\newcommand{\etal}{et al.\@\xspace}
\newcommand{\ie}{i.e.\ }
\newcommand{\eg}{e.g.\ }
\newcommand{\mkkOK}{\checkmark}%\color{green}{\checkmark}
\newcommand{\mkkREQ}{\ding{53}}%requires pifont?%\color{green}{\checkmark}
\newcommand{\mkkNO}{}%\text{\color{red}{\textsf{X}}}

\newlength{\xdescwd}
\makeatletter
\NewEnviron{xdesc}{%
  \setbox0=\vbox{\hbadness=\@M \global\xdescwd=0pt
    \def\item[##1]{%
      \settowidth\@tempdima{\textbf{##1}:}%
      \ifdim\@tempdima>\xdescwd \global\xdescwd=\@tempdima\fi}
  \BODY}
  \begin{description}[leftmargin=\dimexpr\xdescwd+.5em\relax,
    labelindent=0pt,labelsep=.5em,
    labelwidth=\xdescwd,align=left]\BODY\end{description}}
\makeatother

\makeatletter
\newcommand{\boldor}[2]{%
	\ifnum\strcmp{\f@series}{bx}=\z@
		#1%
	\else
		#2%
	\fi
}
\newcommand{\textstyleElProm}[1]{\boldor{\MakeUppercase{#1}}{\textsc{#1}}}
\makeatother
\newcommand{\electre}{\textstyleElProm{Électre}\xspace}
\newcommand{\electreIv}{\textstyleElProm{Électre Iv}\xspace}
\newcommand{\electreIV}{\textstyleElProm{Électre IV}\xspace}
\newcommand{\electreIII}{\textstyleElProm{Électre III}\xspace}
\newcommand{\electreTRI}{\textstyleElProm{Électre Tri}\xspace}
% \newcommand{\utadis}{\texorpdfstring{\textstyleElProm{utadis}\xspace}{UTADIS}}
% \newcommand{\utadisI}{\texorpdfstring{\textstyleElProm{utadis i}\xspace}{UTADIS I}}

%TODO
% \newcommand{\textstyleElProm}[1]{{\rmfamily\textsc{#1}}} 

\newunicodechar{ℕ}{\mathbb{N}}
\newunicodechar{ℝ}{\mathbb{R}}
\newunicodechar{−}{\ifmmode{-}\else\textminus\fi}
\newunicodechar{≠}{\ensuremath{\neq}}
\newunicodechar{≤}{\ensuremath{\leq}}
\newunicodechar{≥}{\ensuremath{\geq}}
\newunicodechar{≻}{\succ}
\newunicodechar{⊁}{\nsucc}
\newunicodechar{⪰}{\succeq}
\newunicodechar{▷}{\triangleright}
\newunicodechar{⋫}{\ntriangleright}
\newunicodechar{→}{\ifmmode\rightarrow\else\textrightarrow\fi}
\newunicodechar{⇒}{\ensuremath{\Rightarrow}}
\newunicodechar{⇏}{\ensuremath{\nRightarrow}}
\newunicodechar{⇔}{\ensuremath{\Leftrightarrow}}
\newunicodechar{∪}{\cup}
\newunicodechar{∩}{\cap}
\newunicodechar{∧}{\land}
\newunicodechar{∨}{\lor}
\newunicodechar{¬}{\ifmmode\lnot\else\textlnot\fi}
\newunicodechar{…}{\ifmmode\ldots\else\textellipsis\fi}
\newunicodechar{×}{\ifmmode\times\else\texttimes\fi}
\newunicodechar{γ}{\ensuremath{\gamma}}
\newunicodechar{□}{\Box}


%const
\newcommand{\tikzboxit}{\path node[draw, overlay, inner sep=0.6mm, fit=(boxed), rectangle] {};}%

\newlength{\GraphsNodeSep}
\setlength{\GraphsNodeSep}{7mm}

% MCDA Drawing Sorting
\newlength{\MCDSCatHeight}
\setlength{\MCDSCatHeight}{6mm}
\newlength{\MCDSAltHeight}
\setlength{\MCDSAltHeight}{4mm}
%separation between two vertical alts
\newlength{\MCDSAltSep}
\setlength{\MCDSAltSep}{2mm}
\newlength{\MCDSCatWidth}
\setlength{\MCDSCatWidth}{3cm}
\newlength{\MCDSEvalRowHeight}
\setlength{\MCDSEvalRowHeight}{6mm}
\newlength{\MCDSAltsToCatsSep}
\setlength{\MCDSAltsToCatsSep}{1.5cm}
\newcounter{MCDSNbAlts}
\newcounter{MCDSNbCats}
\newlength{\MCDSArrowDownOffset}
\setlength{\MCDSArrowDownOffset}{0mm}

\tikzset{/Graphs/dot/.style={
	shape=circle, fill=black, inner sep=0, minimum size=1mm
}}
\tikzset{/MC/D/S/alt/.style={
	shape=rectangle, draw=black, inner sep=0, minimum height=\MCDSAltHeight, minimum width=2.5cm, anchor=north east
}}
\tikzset{MC/D/S/pref/.style={
	shape=ellipse, draw=gray, thick
}}
\tikzset{/MC/D/S/cat/.style={
	shape=rectangle, draw=black, inner sep=0, minimum height=\MCDSCatHeight, minimum width=\MCDSCatWidth, anchor=north west
}}
\tikzset{/MC/D/S/evals matrix/.style={
	matrix, row sep=-\pgflinewidth, column sep=-\pgflinewidth, nodes={shape=rectangle, draw=black, inner sep=0mm, text depth=0.5ex, text height=1em, minimum height=\MCDSEvalRowHeight, minimum width=12mm}, nodes in empty cells, matrix of nodes, inner sep=0mm, outer sep=0mm, row 1/.style={nodes={draw=none, minimum height=0em, text height=, inner ysep=1mm}}
}}

\newlength{\GitCommitSep}
\setlength{\GitCommitSep}{13mm}

\tikzset{/Git/commit/.style={
	shape=rectangle, draw, minimum width=4em, minimum height=0.6cm
}}
\tikzset{/Git/branch/.style={
	shape=ellipse, draw, red
}}
\tikzset{/Git/head/.style={
	shape=ellipse, draw, fill=yellow
}}

\tikzset{profile matrix/.style={
	matrix of math nodes, column sep=3mm, row sep=2mm, nodes={inner sep=0.5mm, anchor=base}
}}
\tikzset{rank-profile matrix/.style={
	matrix of math nodes, column sep=3mm, row sep=2mm, nodes={anchor=base}, column 1/.style={nodes={inner sep=0.5mm}}, row 1/.style={nodes={inner sep=0.5mm}}
}}
\tikzset{rank-vector/.style={
	draw, rectangle, inner sep=0, outer sep=1mm
}}
\tikzset{isolated rank-vector/.style={
	draw, matrix of math nodes, column sep=3mm, inner sep=0, matrix anchor=base, nodes={anchor=base, inner sep=.33em}, ampersand replacement=\&
}}

% GUI
\tikzset{/GUI/button/.style={
	rectangle, very thick, rounded corners, draw=black, fill=black!40%, top color=black!70, bottom color=white
}}

% Logger objects
\tikzset{/logger/main/.style={
	shape=rectangle, draw=black, inner sep=1ex, minimum height=7mm
}}
\tikzset{/logger/helper/.style={
	shape=rectangle, draw=black, dashed, minimum height=7mm
}}
\tikzset{/logger/helper line/.style={
	<->, draw, dotted
}}

% Beliefs
\tikzset{/Beliefs/D/S/attacker/.style={
	shape=rectangle, draw, minimum size=8mm
}}
\tikzset{/Beliefs/D/S/supporter/.style={
	shape=circle, draw
}}

\newcommand{\tikzmark}[1]{%
	\tikz[overlay, remember picture, baseline=(#1.base)] \node (#1) {};%
}


\begin{acronym}
\acro{AMCD}{Aide Multicritère à la Décision}
\acro{ASA}{Argument Strength Assessment}
\acro{DA}{Decision Analysis}
\acro{DM}{Decision Maker}
\acro{DPr}{Deliberated Preferences}
\acro{DRSA}{Dominance-based Rough Set Approach}
\acro{DSS}{Decision Support Systems}
\acrodefplural{DSS}{Decision Support Systems}
% \newacroplural{DSS}[DSSes]{Decision Support Systems}
\acro{EJOR}{European Journal of Operational Research}
\acro{LNCS}{Lecture Notes in Computer Science}
\acro{MCDA}{Multicriteria Decision Aid}
\acro{MIP}{Mixed Integer Program}
\acro{NCSM}{Non Compensatory Sorting Model}
\acro{PL}{Programme Linéaire}
\acro{PLNE}{Programme Linéaire en Nombres Entiers}
\acro{PM}{Programme Mathématique}
\acro{MP}{Mathematical Program}
\acro{MIP}{Mixed Integer Program}
% \newacroplural{PM}{Programmes Mathématiques}
%acrodefplural since version 1.35, my debian has \ProvidesPackage{acronym}[2009/01/25, v1.34, Support for acronyms (Tobias Oetiker)]
\acrodefplural{PM}{Programmes Mathématiques}
\acro{PMML}{Predictive Model Markup Language}
\acro{RESS}{Reliability Engineering \& System Safety}
\acro{SMAA}{Stochastic Multicriteria Acceptability Analysis}
\acro{URPDM}{Uncertainty and Robustness in Planning and Decision Making}
\acro{XML}{Extensible Markup Language}
\end{acronym}

\usepackage{calc}
%\defbeamertemplate*{footline}{default}
%{}

\title[Eliciting a voting rule]{Eliciting a suitable voting rule via examples}
\subject{Social choice theory}
\keywords{preference elicitation, voting rules}
\author{\emph{Olivier Cailloux} \and Ulle Endriss}
\institute[ILLC]{Heudiasyc, University of Technology of Compiègne \& ILLC, University of Amsterdam}
\date{\formatdate{18}{10}{2014}}

\begin{document}
\bibliographystyle{apalike}

\begin{frame}[plain]
	\tikz[remember picture,overlay]{
		\node[anchor=south west, inner sep=0] at (current page.south west) {
			{\includegraphics[width=1cm]{illc_cropped}}
		};
	}
   \titlepage
\end{frame}

\section{Context}
\subsection{Our goal}
\begin{frame}
	\frametitle{Introduction}
	
	\begin{block}{Context}
	\begin{itemize}
		\item A \emph{committee} (a group of decision makers)
		\begin{itemize}
			\item a panel attributing a research price
			\item a management committee
		\end{itemize}
		\item Recurring decisions
		\item A decision is taken using a voting rule
		\item Voting rule: a systematic way of aggregating different opinions and decide
	\end{itemize}
	\end{block}
	\begin{block}{Our goal}
		We want to help the committee choose a suitable voting rule.
	\end{block}
\end{frame}

\begin{frame}[fragile]
	\frametitle{Voting rule}

	\begin{block}{Input}
		\begin{itemize}
			\item A set of possible alternatives (options) $\allalts$
			\item Each voter $i \in N$ has a linear order of preference over $\allalts$
			\item A profile $\prof$ associates each $i$ to such an order.
		\end{itemize}
	\end{block}
	\vfill
	\begin{minipage}[c]{(\linewidth - 3mm) * \real{0.45}}
		\vspace{-5mm}
		\begin{block}{Voting rule}
			Associates to each profile $\prof$ winning alternatives $A \subseteq \allalts$.
		\end{block}
	\end{minipage}\hspace{3mm}%
	\begin{minipage}[c]{(\linewidth - 3mm) * \real{0.55}}
		\mbox{}
		\hfill
		\begin{tikzpicture}
			\path node[profile matrix] (profile) {
				R_1&
				R_2
				\\
				| (profile11) | a&
				b
				\\
				b&
				a
				\\
				c&
				| (profile32) | c
				\\
			};
			\path ($(profile.south west)!.5!(profile.south east)$) ++ (0, -5mm) node {$\prof$};
			
			\path node[draw, rectangle, fit=(profile11) (profile32), outer xsep=2mm, outer ysep=1mm] (justprofile) {};
			\path (justprofile.east) ++ (2.5cm, 0) node[inner sep=0] (winners) {\mbox{} $\Set{a, b}$};
			\path[draw, ->] (justprofile.east) to[bend left=35] node[anchor=south] {$f$} (winners.west);

			\path[draw, decorate, decoration={brace, mirror}] (justprofile.south west) -- (justprofile.south east);
		\end{tikzpicture}
	\end{minipage}
\end{frame}

\begin{frame}
	\frametitle{Our goal}
	
	Making decisions involves two steps.
	\begin{enumerate}
		\item Establish a constitution: choose a voting rule.
		\item Solve a decision problem: apply the voting rule.
	\end{enumerate}
	
	\begin{block}{Our goal}
		We focus on step 1: help the committee choose a voting rule.
		\begin{itemize}
			\item Class of functions $\allrules$ (the set of all voting rules)
			\item Preference elicitation in order to choose a function $f \in \allrules$.
			\item We want to ask \emph{simple} questions: example-based.
		\end{itemize}
	\end{block}
\end{frame}

\subsection{General approach}
\begin{frame}[fragile]
	\frametitle{A naïve attempt}
	
	\begin{block}{A first attempt}
		Simply give a profile $\prof$ and ask for $f(\prof)$. Then iterate.
	\end{block}
\vspace{-4mm}
	\begin{center}
		\begin{tikzpicture}
			\path node[profile matrix] (profile) {
				R_1&
				R_2
				\\
				| (profile11) | a&
				c
				\\
				b&
				b
				\\
				c&
				| (profile32) | a
				\\
			};
			\path ($(profile.south west)!.5!(profile.south east)$) ++ (0, -5mm) node {$\prof$};
			
			\path node[draw, rectangle, fit=(profile11) (profile32), outer xsep=2mm, outer ysep=1mm] (justprofile) {};
			\path (justprofile.east) ++ (2.5cm, 0) node[inner sep=0] (winners) {\mbox{} $\Set{?}$};
			\path[draw, ->] (justprofile.east) to[bend left=35] node[anchor=south] {$f$} (winners.west);

			\path[draw, decorate, decoration={brace, mirror}] (justprofile.south west) -- (justprofile.south east);
		\end{tikzpicture}
	\end{center}\vspace{-2mm}
	\begin{itemize}
		\item Completely general: all functions in $\allrules$ can be reached.
	\end{itemize}
	\emph{But…} 
	\begin{itemize}
		\item One question brings very little information.
		\item Questions may be difficult to answer.
	\end{itemize}
\end{frame}

\begin{frame}
	\frametitle{General idea}
	
	\begin{itemize}
		\item Ask \emph{good} (informative, example-based) questions.
		\item Restrict the class of a-priori acceptable functions to $\allrules' \subset \allrules$.
	\end{itemize}
\end{frame}

\begin{frame}
	\frametitle{Outline}
	\tableofcontents[hideallsubsections]
\end{frame}

\AtBeginSection{
	\begin{frame}
		\frametitle{Outline}
		\tableofcontents[currentsection, hideallsubsections]
	\end{frame}
}

\section{Asking good questions}
\begin{frame}[fragile]
	\frametitle{A different view of a profile}
	
	\begin{itemize}
		\item We want to ask more informative questions about $f$.
		\item We look at profiles under a different angle.
		\item A \emph{rank-vector} maps voters to ranks, $x: N \rightarrow [1, m]$.
		\item All rank vectors: $\rvspace$.
	\end{itemize}
	
	\begin{center}
	\begin{tikzpicture}
			\path node[profile matrix] (profile) {
				R_1&
				R_2
				\\
				| (profile11) | a&
				a
				\\
				b&
				c
				\\
				c&
				| (profile32) | b
				\\
			};
		\path[draw, decorate, decoration={brace, mirror}] (profile.south west) -- (profile.south east);
		\path ($(profile.south west)!.5!(profile.south east)$) ++ (0, -5mm) node {$\prof$};
		
		\path (profile.north east) ++ (1.5cm, 0) node[rank-profile matrix, anchor=north west] (rank-profile) {
			&
			| (header start) | 1&
			| (header end) | 2
			\\
			| (alts start) | a&
			| (rv1 start) | 1&
			| (rv1 end) | 1
			\\
			b&
			| (rv2 start) | 2&
			| (rv2 end) | 3
			\\
			| (alts end) | c&
			| (rv3 start) | 3&
			| (rv3 end) | 2
			\\
		};
		\path node[draw, ellipse, dotted, inner sep=0, fit=(header start.north west) (header end.south east)] (N) {};
		\path (N.north) node[anchor=south, inner sep=1mm] {$N$};
		\path node[draw, ellipse, dotted, inner sep=0, fit=(alts start.north west) (alts end.south east)] (A) {};
		\path (A.west) node[anchor=east, inner sep=1mm] {$\allalts$};
		\foreach \i/\a in {1/a, 2/b, 3/c} {
			\path node[rank-vector, fit=(rv\i\space start.north west) (rv\i\space end.south east)] (rv\i) {};
			\path[<-, draw] (rv\i.east) -- ++(0.3cm, 0) node[anchor=west] {$\rprof_{\prof}(\a) \in [1, m]^N$};
		}
		\path[draw, decorate, decoration={brace, mirror, raise=2mm}] (rank-profile.south west) -- (rank-profile.south east);
		\path ($(rank-profile.south west)!.5!(rank-profile.south east)$) ++ (0, -7mm) node {$\rprof_{\prof}$};
	\end{tikzpicture}
	\end{center}

	To each profile $\prof$ corresponds a \emph{rank-profile} $\rprof_\prof$.
\end{frame}

\subsection{Using a weak-order}
\begin{frame}[fragile]
	\frametitle{Representing the preferences of the committee}
	
	\begin{itemize}
		\item We can now ask for the preference status of, e.g.,\\ \isolatedrankvector{1}{3} versus \isolatedrankvector{2}{2}.
		\item Series of such questions permit to identify a voting rule.
	\end{itemize}
	
	\begin{definition}[Weak-order based rules]
		\begin{itemize}
			\item $\pglobeq$ a weak-order (transitive, reflexive, connected) over $\rvspace$.
			\item Having a profile $\prof$, look at the maximal rank-vectors in it according to $\pglobeq$.
			\item The rule $f_\pglobeq$, at $\prof$, selects those alternatives having maximal rank-vectors in $\prof$.
		\end{itemize}
		A rule $f$ is \emph{weak-order based} if there exists $\pglobeq$ st $f = f_\pglobeq$.
	\end{definition}
\end{frame}
%TODO include an example with layers of rvs to describe the weak-order, and the profile involving 13 and 22, and show who wins.

\subsection{Using preorders}
\begin{frame}
	\frametitle{Incomplete question sets}
	
	\begin{itemize}
		\item We do not want to ask every possible questions!
		\item Can we get away with only \emph{some} answers?
	\end{itemize}
	\begin{definition}[Robust rules]
		\begin{itemize}
			\item $\ppglobeq$ a \emph{preorder} (transitive, reflexive) over $\rvspace$.
			\item Look at all weak-orders $\pglobeq$ extending $\ppglobeq$.
			\item The \emph{robust} rule $F_\ppglobeq$, at $\prof$, selects those alternatives winning in \emph{some} $f_\pglobeq$ (for some $\pglobeq$ extension of $\ppglobeq$).
		\end{itemize}
		A rule $f$ is \emph{robust} if there exists $\ppglobeq$ st $f = F_\ppglobeq$.
	\end{definition}
\end{frame}

\section{Restrict the class of functions}
\subsection{WOB rules}
\begin{frame}
	\frametitle{The WOB class}
	
	Not every rule is weak-order based.
	\begin{itemize}
		\item Bad news: we are not fully general any more.
		\item Good news: we have restricted our class of functions. $WOB = \Set{f_\pglobeq, \pglobeq \text{ a weak-order over } \rvspace}$ instead of $\allrules$.
	\end{itemize}
	How does WOB compare to other known classes of rules?
	\begin{itemize}
		\item Every scoring rule (e.g. Borda) is a WOB rule.
		\item Some WOB rules are not scoring rules. E.g. Bucklin.
		\item Many Condorcet rules are not WOB rules.
	\end{itemize}
	\begin{block}{Some relationships between classes of rules}
		Scoring $\subset$ WOB; Condorcet $\cap$ WOB $= \emptyset$ (for $n = 3k, m ≥ 4$).
	\end{block}
 \end{frame}

\subsection{Robust rules}
\begin{frame}[fragile]
	\frametitle{The class of robust rules}
	
	Some robust rules are not WOB rules.
	\begin{example}[A robust rule]
		\begin{minipage}{0.45\columnwidth}
			\centering
			${\ppglobeq} = {}$\fbox{%
				\begin{tikzpicture}[baseline=(current bounding box)]
\setlength{\GraphsNodeSep}{11mm}
					\path node[isolated rank-vector, matrix anchor=east, outer sep=0, inner sep=0] (13) {1\&3\\};
					\path (13.south) ++(0, -1\GraphsNodeSep) node[isolated rank-vector, matrix anchor=north] (41) {4\&1\\};
					\path (41.south) ++(-8mm, -1\GraphsNodeSep) node[isolated rank-vector, matrix anchor=north, inner sep=0] (14) {1\&4\\};
					\path (41.south) ++(8mm, -1\GraphsNodeSep) node[isolated rank-vector, matrix anchor=north, inner sep=0] (33) {3\&3\\};
					\path (33.east) ++(4mm, 0) node[isolated rank-vector, matrix anchor=west, inner sep=0] (34) {3\&4\\};
					\path ($(13 -| 34)!.5!(41 -| 34)$) node[isolated rank-vector, matrix anchor=center, inner sep=0] (22) {2\&2\\};
					\path[->, draw] (13) -- (41);
					\path[->, draw] (41) -- (14);
					\path[->, draw] (41) -- (33);
		%			\path[->, draw] (41) -- (34);
					\path[->, draw] (22) -- (34);
				\end{tikzpicture}%
			}
		\end{minipage}%
		\begin{minipage}{0.55\columnwidth}
			\centering
			$\begin{array}{lcl@{\hspace*{5mm}}cl}
				& \multicolumn{2}{c}{\hspace*{3mm}\prof_1\hspace{3mm}F_\ppglobeq(\prof_1)}	& \multicolumn{2}{c}{\hspace*{2mm}\prof_2\hspace{3mm}F_\ppglobeq(\prof_2)}\\
				\cmidrule(r{2mm}){2-3}\cmidrule{4-5}
%		\begin{tikzpicture}[baseline]
%			\path node[draw, inner ysep=1cm, outer sep=0, anchor=base] (first) {1};
%			\path (current bounding box.south west) ++(0, -1cm);
%		\end{tikzpicture}%
				a &\isolatedrankvectorinarray{1}{3}
 & \mkkOK & \isolatedrankvectorinarray{1}{4} &\\
				b &\isolatedrankvectorinarray{2}{2} & \mkkOK & \isolatedrankvectorinarray{2}{2} & \mkkOK\\
				c &\isolatedrankvectorinarray{3}{4} & & \isolatedrankvectorinarray{3}{3} & \\
				d &\isolatedrankvectorinarray{4}{1} & & \isolatedrankvectorinarray{4}{1} & \mkkOK\\
			\end{array}$
		\end{minipage}
	\end{example}
	
	\begin{block}{Some relationships between classes of rules}
		Scoring $\subset$ WOB $\subset$ Robust.
	\end{block}
\end{frame}

\section{Which questions to ask?}
\begin{frame}
	\frametitle{Which questions to ask?}
	
	\begin{itemize}
		\item Assume the committee has a weak-order $\pglobeq$ in mind.
		\item We want to discover much information using few questions.
		\item Different questions bring different amount of information.
	\end{itemize}
	\begin{definition}[Elicitation strategy]
		An elicitation strategy tells us which question should be asked considering our current knowledge.
	\end{definition}
	A strategy:
	\begin{enumerate}
		\item computes the fitness of asking about a pair of rank-vectors, for each pair;
		\item chooses the fittest pair.
	\end{enumerate}
	We ask $q$ questions, then compare our approximation $F_\ppglobeq$ to $f_\pglobeq$.
\end{frame}

\subsection{Comparing strategies}
\begin{frame}
	\frametitle{Which strategy?}
	
	We tested three strategies.
	\begin{description}
		\item[optimistic] fitness of ($x, y$) proportional to the number of rank-vectors dominated by $x$ or $y$, but not both;
		\item[pessimistic] a variant of the previous strategy, using the min operator rather than the sum;
		\item[likelihood] fitness proportional to the likelihood of a profile occurring where both rank-vectors are possible winners.% (depends on the probability distribution over profiles, we assumed impartial culture)
	\end{description}
	We assume pareto-dominance and indifference to permutations.
	\begin{block}{Comparison of strategies}
		\begin{itemize}
			\item Optimistic not better than random!
			\item Likelihood much better than pessimistic.
		\end{itemize}
	\end{block}
\end{frame}

\subsection{Number of questions}
\begin{frame}
	\frametitle{Number of questions}
	
	How many questions must be asked for a useful approximation?
	\begin{itemize}
		\item Our approximation has all the true winners: $f_\pglobeq(\prof) \subseteq F_\ppglobeq(\prof)$.
		\item But it may have supplementary winners.
		\item We are interested in the ratio of approximated VS true winners: $\frac{\card{F_\ppglobeq(\prof)}}{\card{f_\pglobeq(\prof)}}$.
		\item We average it over all profiles: $\frac{1}{\card{\profiles}} \sum_{\prof \in \profiles} \frac{\card{F_\ppglobeq(\prof)}}{\card{f_\pglobeq(\prof)}}$.
	\end{itemize}
	
	For $6$ voters, $6$ alternatives, using the likelihood strategy:
	\begin{center}
		\begin{tabular}{rrr}
			\toprule
			 & \multicolumn{2}{c}{Target rule}\\
			nb q & Borda & Random $\pglobeq$\\
			\midrule
			0 & 1.9 & 2.2\\
			25 & 1.3 & 1.7\\
			99 & 1.0 & 1.3\\
			\bottomrule
		\end{tabular}
	\end{center}
\end{frame}

\section{Conclusion}
\begin{frame}
	\frametitle{Conclusion}
	
	We propose to help a committee choose a voting rule.
	\begin{itemize}
		\item We introduce a different look at a profile \citep[see also][]{sen_weights_1977}.
		\item We use it to ask simple questions to elicit preferences.
		\item We analyse the class of rules reachable by our questioning process.
		\item A robust voting rule may be defined to give all possible winners \citep[inspired by][]{dias_aggregation/disaggregation_2002}.
		\item We compare and analyse several elicitation strategies.
	\end{itemize}
\end{frame}

\begin{frame}[plain]
	\addtocounter{framenumber}{-1}
	\begin{center}
		\huge
		\textit{Thank you for your attention!}
	\end{center}
	\tikz[remember picture,overlay]
	\node[anchor=south west, inner sep=0] at (current page.south west) {
		{\includegraphics[width=1cm]{illc_cropped}}
	};
\end{frame}

\appendix
\AtBeginSection{
}

\begin{frame}%[allowframebreaks]
	\frametitle{Bibliography}
	\def\newblock{\hskip .11em plus .33em minus .07em}
 	\bibliography{zotero}
\end{frame}

\end{document}
